% Part: methods
% Chapter: induction
% Section: induction-on-N

\documentclass[../../../include/open-logic-section]{subfiles}

\begin{document}

\olfileid{mth}{ind}{inN}

\olsection{استقرا بر~$\Nat$}

در ساده‌ترین شکل، استقرا فنی برای اثبات نتایج دربارهٔ همهٔ اعداد
طبیعی است. این فن از این واقعیت استفاده می‌کند که با آغاز از $0$ و
افزودن پیاپی~$1$، سرانجام به هر عدد طبیعی می‌رسیم. پس برای اثبات اینکه
چیزی دربارهٔ هر عدد صادق است، می‌توانیم (۱)~نشان دهیم دربارهٔ $0$ صادق
است و (۲)~نشان دهیم هرگاه دربارهٔ عدد~$n$ صادق باشد، دربارهٔ عدد بعدی
~$n+1$ نیز صادق است. اگر عبارت «عدد~$n$ خاصیت $P$ را دارد» با $P(n)$
و «عدد~$k$ خاصیت $P$ را دارد» با $P(k)$ و مانند آن کوتاه شود، اثبات
استقرایی $P(n)$ برای همهٔ $n \in \Nat$ از این دو بخش تشکیل می‌شود:
\begin{enumerate}
\item اثبات $P(0)$، و
\item اثبات اینکه برای هر~$k$، اگر $P(k)$ آنگاه $P(k+1)$.
\end{enumerate}
برای آنکه موضوع کاملاً روشن شود، فرض کنید هم (۱) و هم (۲) را داریم.
(۱) می‌گوید $P(0)$ صادق است. با داشتن (۲)، به‌ویژه می‌دانیم اگر
$P(0)$ آنگاه $P(0+1)$، یعنی $P(1)$. این نتیجه با گذاشتن~$0$ به جای
~$k$ در گزارهٔ کلی «برای هر~$k$، اگر $P(k)$ آنگاه $P(k+1)$» به دست
می‌آید. پس با وضع مقدم، $P(1)$ را داریم. بار دیگر از (۲)، این بار با
گذاشتن~$1$ به جای~$k$، داریم: اگر $P(1)$ آنگاه $P(2)$. چون اکنون
$P(1)$ را اثبات کرده‌ایم، با وضع مقدم $P(2)$ را به دست می‌آوریم، و
همین‌طور ادامه می‌دهیم. برای هر عدد~$n$، پس از $n$ بار انجام این کار،
سرانجام به~$P(n)$ می‌رسیم. پس (۱) و (۲) با هم $P(n)$ را برای هر
$n \in \Nat$ برقرار می‌کنند.

مثالی را بررسی کنیم. فرض کنید می‌خواهیم بدانیم با $n$ تاس چند مجموع
متفاوت می‌توان آورد. هرچند شاید عجیب به نظر برسد، از $0$ تاس آغاز
کنیم. اگر هیچ تاسی نداشته باشید، فقط یک مجموع ممکن است: هیچ خالی
وجود ندارد و مجموع برابر~$0$ است. پس تعداد حالت‌های ممکن~$1$ است. اگر
فقط یک تاس داشته باشید، یعنی $n=1$، شش مقدار ممکن از $1$ تا~$6$ وجود
دارد. با دو تاس هر مجموع از $2$ تا $12$ را می‌توان آورد؛ یعنی $11$
حالت. با سه تاس هر عدد از $3$ تا $18$ را می‌توان آورد؛ یعنی $16$
حالت متفاوت. دنبالهٔ $1$، $6$، $11$، $16$ الگویی را نشان می‌دهد:
شاید پاسخ $5n+1$ باشد؟ البته $5n+1$ بیشترین تعداد ممکن است، چون میان
$n$، کمترین مقدار ممکن با $n$ تاس که وقتی همه $1$ باشند حاصل می‌شود،
و $6n$، بیشترین مقدار ممکن که وقتی همه $6$ باشند حاصل می‌شود، فقط
$5n+1$ عدد وجود دارد.

\begin{thm}
  با $n$ تاس می‌توان همهٔ $5n+1$ مقدار ممکن میان $n$ و $6n$ را آورد.
\end{thm}

\begin{proof}
بگذارید $P(n)$ این گزاره باشد: «با $n$ تاس می‌توان هر عدد میان $n$
و~$6n$ را آورد.» برای استفاده از استقرا، دو بخش زیر را اثبات می‌کنیم:
\begin{enumerate}
\item \emph{پایهٔ استقرا}، یعنی $P(0)$: با صفر تاس فقط مجموع $0$ ممکن
  است، و این تنها عدد میان $0$ و~$0$ است.
\item \emph{گام استقرا}: برای هر $k$، اگر $P(k)$ آنگاه~$P(k+1)$.
\end{enumerate}

(۱) بی‌درنگ برقرار است: با صفر تاس هیچ خالی نمی‌آید و تنها مجموع
ممکن~$0$ است. پس با صفر تاس می‌توان هر عدد میان $0$ و~$0$ را آورد.

برای اثبات (۲)، مقدم شرط، یعنی $P(k)$، را فرض می‌کنیم. این فرض را
\emph{فرض استقرا} می‌نامند. از آن برای اثبات~$P(k+1)$ استفاده می‌کنیم.
بخش دشوار این است که راهی بیابیم تا مقادیر ممکن پرتاب $k+1$ تاس را
بر حسب مقادیر ممکن پرتاب $k$ تاس و پرتاب تاس اضافیِ $(k+1)$ام در نظر
بگیریم؛ بااین‌حال، اگر بخواهیم فرض استقرا را به کار ببریم باید همین
کار را انجام دهیم.

فرض استقرا می‌گوید با $k$ تاس می‌توان هر عدد میان $k$ و~$6k$ را آورد.
اگر با تاس $(k+1)$ام عدد~$1$ بیاوریم، $1$ به مجموع افزوده می‌شود. پس
با پرتاب $k$ تاس و سپس آوردن~$1$ با تاس $(k+1)$ام، می‌توان هر مقدار
میان $k+1$ و $6k+1$ را آورد. چه چیز باقی می‌ماند؟ مقادیر $6k+2$ تا
$6k+6$. این مقادیر با آوردن $6$ روی هر یک از $k$ تاس نخست و سپس عددی
میان $2$ و $6$ روی تاس $(k+1)$ام حاصل می‌شوند. در مجموع، این یعنی با
$k+1$ تاس می‌توان هر عدد میان $k+1$ و $6(k+1)$ را آورد؛ یعنی با فرض
استقرا~$P(k)$، گزارهٔ~$P(k+1)$ را اثبات کرده‌ایم.
\end{proof}

بسیار پیش می‌آید که برای اثبات چیزی دربارهٔ دنباله‌ای از اشیا — مانند
اعداد یا مجموعه‌ها — از استقرا استفاده کنیم، وقتی خود دنباله «به‌طور
استقرایی» تعریف شده باشد؛ یعنی شیء $(n+1)$ام بر حسب شیء $n$ام تعریف
شود. برای نمونه، می‌توانیم مجموع~$s_n$ اعداد طبیعی تا~$n$ را چنین
تعریف کنیم:
\begin{align*}
  s_0 & = 0\\
  s_{n+1} & = s_n + (n+1)
\end{align*}
این تعریف می‌دهد:
\begin{align*}
  s_0 & = 0,\\
  s_1 & = s_0 + 1 && = 1,\\
  s_2 & = s_1 + 2 && = 1 + 2 = 3\\
  s_3 & = s_2 + 3 && = 1 + 2 + 3 = 6, \text{ و مانند آن.}
\end{align*}
اکنون می‌توانیم با استقرا اثبات کنیم $s_n = n(n+1)/2$.

\begin{prop}
  $s_n = n(n+1)/2$.
\end{prop}

\begin{proof}
  باید اثبات کنیم (۱)~$s_0 = 0\cdot(0 + 1)/2$ و (۲)~اگر
  $s_k = k(k+1)/2$، آنگاه $s_{k+1} = (k+1)(k+2)/2$. بخش (۱) آشکار
  است. برای اثبات (۲)، فرض استقرا را می‌پذیریم:
  $s_k = k(k+1)/2$. با استفاده از آن باید نشان دهیم
  $s_{k+1} = (k+1)(k+2)/2$.

  مقدار $s_{k+1}$ چیست؟ بنا بر تعریف، $s_{k+1} = s_k + (k+1)$. بنا بر
  فرض استقرا، $s_k = k(k+1)/2$. می‌توانیم این مقدار را در معادلهٔ پیش
  جای‌گذاری کنیم و سپس فقط اندکی حساب کسرها لازم است:
  \begin{align*}
    s_{k+1} & = \frac{k(k+1)}{2} + (k+1) \\
    & = \frac{k(k+1)}{2} + \frac{2(k+1)}{2} \\
    & = \frac{k(k+1) + 2(k+1)}{2} \\
    & = \frac{(k+2)(k+1)}{2}.
  \end{align*}
\end{proof}

درس مهم این است که وقتی چیزی را دربارهٔ دنبالهٔ به‌طور استقرایی
تعریف‌شده‌ای مانند $a_n$ اثبات می‌کنید، استقرا راه آشکار پیش روست.
حتی اگر چنین نباشد، مانند مسئلهٔ حالت‌های پرتاب تاس، هرگاه بتوانید
حالت~$k+1$ را به حالت~$k$ مربوط کنید می‌توانید از استقرا بهره بگیرید.

\end{document}
